% ============================================================
% 分片：s08-functions  对应原文 PDF 第 96--111 页（印刷页 96--111）
% 负责章节：§8 From Hecke Eigensheaves To Hecke Eigenfunctions + 8.1--8.10
% 译者：trans-s08   审校状态：初稿
% ============================================================

% ==== 原文 PDF 第 96 页（印刷页 96）====
% [跨页接续: 原文第 96 页开头四行仍属 §7（由 s07 分片负责），本分片自该页 §8 标题起]

\section{从 Hecke 本征层到 Hecke 本征函数}

\begin{keytakeaway}
\textbf{核心物理图像：}经典一侧的 Hecke 本征函数与几何一侧的 Hecke 本征层由 Grothendieck 的 faisceaux--fonctions 词典联系：在有限域上，层经「取 Frobenius 在茎上的迹」变为函数，而层上的 Hecke 函子变为函数上的 Hecke 算子。问题是：$\mathbb{F}_q$ 上的「朴素」本征层性质只保证在 $C$ 的 $\mathbb{F}_q$-点处的本征性；要让本征性对 $C$ 的所有闭点成立，必须对 Hecke 函子补加 $S_2$（一般 $S_m$）等变性，即 Hecke 函子的交换性条件——这正是 Tannakian 意义下纤维函子所需的最后一个约束。
\textbf{数学工具：}$\ell$-进反常层与导出范畴、几何 Frobenius 与 Frobenius 迹、Hecke 函子的迭代与对称化、几何 Satake 等价、交换类域论（$\mathrm{Pic}(\mathbb{F}_q)$ 与 $W_F$ 的非分歧阿贝尔商）、$SL_2$ 的 Labesse--Langlands 型全局 L-包判据（定理 7.1）、对有限群 $\mathbb{Z}_2$ 的 Fourier 变换。
\textbf{主要结论：}内窥 $PGL_2$-局部系统（像约化到 $O_2$）所对应的正则 Hecke 本征层分裂为两个分数 Hecke 本征层 $\mathcal{F}_\pm$；其函数 $f_+ + f_-$ 与 $f_+ - f_-$ 分别是对 $\sigma$ 与 $\sigma' = \sigma\otimes\nu$ 的通常 Hecke 本征函数，而 $f_\pm$ 本身在奇数次闭点处满足一套本征值、在偶数次闭点处满足另一套（(8.19)、(8.22)）。improper Hecke 函子把两个不可约对象都送到同一对象，故不是范畴等价，这正是经典中 improper Hecke 算子消失、L-包结构非平凡的几何对应物；而在 Iwahori 情形（(8.28)）中这一消失现象不再出现。
\end{keytakeaway}

现在我们希望用 Hecke 本征层替换 Hecke 本征函数：这些几何对象使我们能够把经典 Langlands 对应与几何 Langlands 对应联系起来，并最终与前面各节所讨论的对偶群的 Hitchin 纤维化的镜像对称联系起来。

\subsection{特征 \texorpdfstring{$p$}{p} 中的 Hecke 本征层}

在 6.3 节中先前对 Hecke 本征层的讨论里，我们曾假定曲线定义在 $\mathbb{C}$ 上。此时，对应于 $C$ 上 ${}^{L}G$-局部系统 $\mathcal{E}$ 的 Hecke 本征层，是 $\Bun_G$ 上的一个 $\CD$-模 $\mathcal{F}$，连同同构 $\alpha_V$ 作为附加数据（见公式 (6.5)）。\footnote{预期该 $\CD$-模是 holonomic 的，并具有正则奇点。}借助 Riemann--Hilbert 对应，我们便可从 $\CD$-模切换到反常层（perverse sheaf）（例如见 [F1] 第 3.4 节）。于是，Hecke 本征层可以看作 $\Bun_G$ 上反常层范畴的对象，并配有同构 (6.5)。

% [Tutor Note: 跳步补足。原文「借助 Riemann--Hilbert 对应从 D-模切换到反常层」跳过了一个关键前提：Riemann--Hilbert 对应只在正则奇性的 holonomic D-模上与反常层之间才是范畴等价（Deligne 的版本须在正则 holonomic 情形使用同一 t-结构）；这正是脚注 43 要强调「holonomic 且正则奇点」的原因。若奇异点是正则的，$\mathcal{D}$-模的解局部系统给出反常层；若奇异点不规则，则须改用 Stokes 结构（扭 D-模/指数 D-模）。本文 §6 起所有 Hecke 本征层都由 opers 的 Hecke 特征构造，均属正则 holonomic 情形，故这一步是合法的。 来源：原书 p.96 上下文（A）。]

现在我们把复曲线换成定义在有限域 $k = \mathbb{F}_q$ 上的曲线 $C$。特征 $0$ 中反常层的概念，对有限域上的代数簇（或代数叠）有对应物（它们是 $\ell$-进层的导出范畴的对象 [BBD]）。我们有曲线 $C$ 上 $G$-丛的模叠 $\Bun_G$，其中 $C$ 定义在 $k$ 上。这是 $k$ 上的代数叠。因此我们就有 $\Bun_G$ 上对应于非分歧同态 $\sigma : W_F \to {}^{L}G$ 的 Hecke 本征层的概念。也就是说，我们把 $\sigma$ 视为 $C$ 上的 $\ell$-进 ${}^{L}G$-局部系统 $\mathcal{E}$。换言之，对 ${}^{L}G$ 的每个表示 $V$，相应的扭转
\[
\mathcal{V}_{\mathcal{E}} = \mathcal{E} \times_{{}^{L}G} V
\]
是 $C$ 上局部常数的 $\ell$-进层，并且这些层与 ${}^{L}G$ 的表示上的张量积结构相容。我们还有 Hecke 函子 $H_V$，$V\in \mathrm{Rep}({}^{L}G)$，其定义方式与 $\mathbb{C}$ 上相同。

按定义，以 $\mathcal{E}$（或 $\sigma$）为「本征值」的 Hecke 本征层，是 $\Bun_G$ 上的一个反常（$\ell$-进）层 $\mathcal{F}$，连同同构作为附加数据（与 (6.5) 比较）
\begin{equation}\tag{8.1}
\alpha_V : H_V(\mathcal{F}) \xrightarrow{\ \sim\ } \mathcal{V}_{\mathcal{E}} \boxtimes \mathcal{F}.
\end{equation}
这些同构应与两侧的张量积结构以及结合性相容。下文将看到，为保证从 Hecke 本征层到 Hecke 本征函数的过渡，我们需要再施加一个等变性条件。

为更详细地解释这一过渡，我们回忆：对 $\mathbb{F}_q$ 上的任意代数簇（或代数叠）$Y$，我们可以给 $Y$ 上的任意 $\ell$-进层（或复形）$\mathcal{F}$ 指定 $Y$ 的 $\mathbb{F}_q$-点之集上的一个函数（见 [De, La1]）。事实上，设 $y$ 是 $Y$ 的一个
% ==== 原文 PDF 第 97 页（印刷页 97）====
$\mathbb{F}_q$-点，而 $\bar y$ 是对应于嵌入 $\mathbb{F}_q \hookrightarrow \overline{\mathbb{F}}_q$ 的 $\overline{\mathbb{F}}_q$-点。则 $\mathcal{F}$ 沿复合 $\bar y \to y \to Y$ 的拉回是点 $\operatorname{Spec}\overline{\mathbb{F}}_q$ 上的一个（$\ell$-进）层。这样一个层的数据等同于一个 $\overline{\mathbb{Q}}_\ell$-向量空间的数据，我们可以把它看作 $\mathcal{F}$ 在 $\bar y$ 处的茎 $\mathcal{F}_{\bar y}$。这个向量空间上还有一份附加数据。事实上，Galois 群 $\mathrm{Gal}(\overline{\mathbb{F}}_q/\mathbb{F}_q)$ 是态射 $\bar y \to y$ 的对称群，因此它作用在 $\mathcal{F}_{\bar y}$ 上。特别地，我们有（几何）Frobenius 元素 $\mathrm{Fr}_y$ 的作用，它对应于 $\mathbb{F}_q$ 的 Galois 群的生成元（的逆），作用为 $x \mapsto x^q$。这个自同构依赖于态射 $\bar y \to y$ 的选取，但它的共轭类不依赖于任何选取。于是我们得到茎 $\mathcal{F}_{\bar y}$ 的一个自同构的共轭类。因此，几何 Frobenius 自同构的迹可以典范地指定给 $\mathcal{F}$ 与 $y$。我们用 $\mathrm{Tr}(\mathrm{Fr}_y, \mathcal{F})$ 表示它。

更一般地，若 $\mathcal{F}$ 是 $\ell$-进层的复形，我们取 $\mathrm{Fr}_y$ 在 $\mathcal{F}$ 于 $\bar y$ 处的茎上同调上的迹的交错和。于是我们得到 $Y$ 的 $\mathbb{F}_q$-点之集上的一个函数 $\mathrm{f}_{\mathcal{F},\mathbb{F}_q}$，它在 $y\in Y(\mathbb{F}_q)$ 处的值为
\[
\mathrm{f}_{\mathcal{F},\mathbb{F}_q}(y) = \sum_i (-1)^i\,\mathrm{Tr}\bigl(\mathrm{Fr}_y, H^i_{\bar y}(\mathcal{F})\bigr).
\]
类似地，对每个 $n > 1$，我们通过公式定义 $Y$ 的 $\mathbb{F}_{q^n}$-点之集上的函数 $\mathrm{f}_{\mathcal{F},\mathbb{F}_{q^n}}$
\[
\mathrm{f}_{\mathcal{F},\mathbb{F}_{q^n}}(y) = \sum_i (-1)^i\,\mathrm{Tr}\bigl(\mathrm{Fr}_y, H^i_{\bar y}(\mathcal{F})\bigr), \qquad y\in Y(\mathbb{F}_{q^n})
\]
（此时 $\mathrm{Fr}_y$ 对应于自同构 $y \mapsto y^{q^n}$）。

映射 $\mathcal{F} \mapsto \mathrm{f}_{\mathcal{F},\mathbb{F}_{q^n}}$ 把层复形上的自然运算与函数上的自然运算相互交织（见 [La1] 第 1.2 节）。例如，层的拉回对应于函数的拉回，而紧支集层的前推对应于函数的逐纤维积分。这一从层到函数的过渡被称为 Grothendieck 的 \textit{faisceaux--fonctions} 词典。

若 $Y = \Bun_G$，则 $Y$ 的 $\mathbb{F}_q$-点之集自然同构于双商
\begin{equation}\tag{8.2}
G(F)\backslash G(\mathbb{A}_F)/G(\mathcal{O}_F),
\end{equation}
其中
\[
\mathcal{O}_F = \prod_{x\in C}\mathcal{O}_x
\]
（例如见 [F1] 第 3.2 节）。因此 $\Bun_G$ 上任一反常层 $\mathcal{F}$ 都给出双商 (8.2) 上的函数 $\mathrm{f}_{\mathcal{F},\mathbb{F}_q}$。现在设 $(\mathcal{F}, (\alpha_V))$ 是 $\Bun_G$ 上的一个 Hecke 本征层。考虑 $\Bun_G(\mathbb{F}_q)$（同构于双商 (8.2)）上相应的函数 $\mathrm{f}_{\mathcal{F},\mathbb{F}_q}$，以及它在 Hecke 函子 $H_V$ 下的变换，限制在
\[
(C \times \Bun_G)(\mathbb{F}_q) = C(\mathbb{F}_q)\times \Bun_G(\mathbb{F}_q).
\]
Hecke 函子 $H_V$ 在层上的作用变成相应的 Hecke 算子 $T_{V,x}$ 在函数上的作用。因此对每个 $x\in C(\mathbb{F}_q)$，(8.1) 的左端
% ==== 原文 PDF 第 98 页（印刷页 98）====
给出函数 $T_{V,x}\cdot \mathrm{f}_{\mathcal{F},\mathbb{F}_q}$，而右端成为 $\mathrm{Tr}(\mathrm{Fr}_x, \mathcal{V}_{\mathcal{E}})\mathrm{f}_{\mathcal{F},\mathbb{F}_q}$。于是同构 (8.1) 蕴含
\begin{equation}\tag{8.3}
T_{V,x}\cdot \mathrm{f}_{\mathcal{F},\mathbb{F}_q} = \mathrm{Tr}(\mathrm{Fr}_x, \mathcal{V}_{\mathcal{E}})\,\mathrm{f}_{\mathcal{F},\mathbb{F}_q} = \mathrm{Tr}(\sigma_x(\mathrm{Fr}_x), V)\,\mathrm{f}_{\mathcal{F},\mathbb{F}_q}, \qquad \forall x\in C(\mathbb{F}_q)
\end{equation}
（更多细节见 [F1] 第 3.8 节）。

% [Tutor Note: 推导补足。(8.3) 的右端用了两步被略去的推理。其一，Hecke 函子 $H_V$ 在 $x$ 处的作用是把 $\mathcal{F}$ 经 Hecke 对应 $_{x}\mathcal{H}\subset C\times\Bun_G$ 拉回、沿 $x$ 方向作「修改」再前推；在层层面这是拉回--前推的复形，在函数层面就是 Hecke 算子的积分核作用，即 $T_{V,x}\cdot \mathrm{f} = \int_{_{x}\mathcal{H}_y}\mathrm{f}$。其二，把 (8.1) 的右端 $\mathcal{V}_{\mathcal{E}}\boxtimes\mathcal{F}$ 限制到 $x\times\Bun_G$ 时，$\mathcal{V}_{\mathcal{E}}$ 的茎是 $\ell$-进局部系统的纤维，Frobenius 在其上的迹为 $\mathrm{Tr}(\mathrm{Fr}_x,\mathcal{V}_{\mathcal{E}})$，而后者等于 $\mathrm{Tr}(\sigma_x(\mathrm{Fr}_x),V)$——这里 $\sigma_x$ 是 $\sigma$ 在分解群 $W_{F_x}\subset W_F$ 上的限制，而 $V$ 上的迹与用 ${}^{L}G$-局部系统扭转再取纤维的迹一致，因为 $\mathcal{V}_{\mathcal{E}}$ 就是 $\mathcal{E}$ 经 $V$ 的伴随丛。注意 $\mathrm{Tr}(\mathrm{Fr}_x,\cdot)$ 内的第二个变量是层，$\mathrm{Tr}(\sigma_x(\mathrm{Fr}_x),\cdot)$ 内的第二个变量是表示，两者不可混写。 来源：原书 p.98 上下文（A）。]

这就是我们所寻求的 Hecke 本征函数性质，但有一个警示：先验地，这一条件只对 $C$ 的 $\mathbb{F}_q$-点成立。相反，关于 $\sigma$ 的非分歧 Hecke 本征函数，应当对 $C$ 的所有闭点（其剩余域任意）都是 Hecke 算子的本征函数。为保证这一性质对函数 $\mathrm{f}_{\mathcal{F},\mathbb{F}_q}$ 在所有点 $x\in C$ 处成立，我们必须对反常层 $\mathcal{F}$ 施加一个附加条件；即 [FGV] 第 1.1 节中迭代 Hecke 函子的 $S_2$-等变性。这将在下一节讨论。

\subsection{Hecke 本征层的等变性与交换性条件}

回忆 Hecke 函子 $H_V$ 从 $\Bun_G$ 上层之导出范畴作用到 $C\times \Bun_G$ 上层之导出范畴。再次应用这个函子，我们得到迭代 Hecke 函子 $H_V^{\boxtimes 2}$，从 $\Bun_G$ 上层之导出范畴作用到 $C\times C\times \Bun_G$ 上层之导出范畴。以 $\mathcal{E}$ 为「本征值」的 Hecke 本征层 $\mathcal{F}$ 配有同构
\[
\alpha_V : H_V(\mathcal{F}) \simeq \mathcal{V}_{\mathcal{E}} \boxtimes \mathcal{F},
\]
它给出同构
\[
\alpha_V^{\boxtimes 2} : H_V^{\boxtimes 2}(\mathcal{F}) \simeq \mathcal{V}_{\mathcal{E}} \boxtimes \mathcal{V}_{\mathcal{E}} \boxtimes \mathcal{F}.
\]
在对角线 $\Delta\subset C\times C$ 之外，这个同构两侧都有对称群 $S_2$ 的自然作用。我们需要施加的额外条件是：$\alpha_V^{\boxtimes 2}$ 是 $S_2$-等变同构。

% [Tutor Note: 易错预警。这里「离开对角线 $\Delta$」是全部技术要点：在 $C\times C$ 的对角线上两个 Hecke 修改发生在同一点，Hecke 对应退化，$S_2$ 的作用没有定义；只有限制到 $C\times C\setminus\Delta$（即两个修改点不同）时，交换两个因子的换位才给出 $S_2$-作用。同理，后文 $C^m$ 上的 $S_m$-等变性只在去掉「成对对角线之并 $\Delta$」后才陈述。这也是 (8.4) 中 $p:C^m\to\mathrm{Sym}^m C$ 必须避开 $\Delta$ 才成为平展（$S_m$-主）覆盖的原因。 来源：原书 p.98 上下文（A）。]

这一条件蕴含：对从 $\Bun_G$ 上层之导出范畴作用到 $C^m\times \Bun_G$ 上层之导出范畴的第 $m$ 次迭代 Hecke 函子 $H_V^{\boxtimes m}$，同构
\[
\alpha_V^{\boxtimes m} : H_V^{\boxtimes m}(\mathcal{F}) \simeq (\mathcal{V}_{\mathcal{E}})^{\boxtimes m}\boxtimes \mathcal{F}
\]
在 $C^m$ 中成对对角线之并 $\Delta$ 之外是 $S_m$-等变的。

% [原文疑误: 原文作 "the mth iterated Hecke functor $H_V^{\boxtimes m}$ acting from from the derived category"，重复了一个 from；译文按单次 from 处理，不影响数学内容。 来源：原书 p.98 英文原文与译本对照（A）。]

设 $S_2$-等变条件成立。则 $\mathcal{F}$ 是关于对称化 Hecke 函子 $H_V^{(m)}$ 的本征层，后者从 $\Bun_G$ 上层之导出范畴作用到 $(\mathrm{Sym}^m C\setminus\Delta)\times \Bun_G$ 上层之导出范畴；也就是说，在 $(\mathrm{Sym}^m C\setminus\Delta)\times \Bun_G$ 上有同构
\[
\alpha_V^{(m)} : H_V^{(m)}(\mathcal{F}) \simeq \mathcal{V}_{\mathcal{E}}^{(m)}\boxtimes \mathcal{F},
\]
其中
\begin{equation}\tag{8.4}
\mathcal{V}_{\mathcal{E}}^{(m)} = \bigl(p_*\bigl((\mathcal{V}_{\mathcal{E}})^{\boxtimes m}\bigr)\bigr)^{S_m},
\end{equation}
而 $p : C^m \to \mathrm{Sym}^m C$ 是对称化映射。

现在注意，$C$ 的任一 $m$ 次闭点 $x$ 都给出 $\mathrm{Sym}^m C$ 中的一个 $\mathbb{F}_q$-点 $D(x)$（一个 $m$ 次有效除子）。此外，容易看出
\[
\mathrm{Tr}(\mathrm{Fr}_x, \mathcal{V}_{\mathcal{E}}) = \mathrm{Tr}\bigl(\mathrm{Fr}_{D(x)}, \mathcal{V}_{\mathcal{E}}^{(m)}\bigr).
\]
% ==== 原文 PDF 第 99 页（印刷页 99）====
因此，把 $\alpha_V^{(m)}(\mathcal{F})$ 限制到 $D(x)\times \Bun_G$ 并在那里的 $\mathbb{F}_q$-点上计算 Frobenius 的迹，我们发现公式 (8.3) 对 $m$ 次的所有闭点 $x$ 成立。于是，$S_2$-等变条件保证函数 $\mathrm{f}_{\mathcal{F},\mathbb{F}_q}$ 确实是 (8.2) 上关于局部系统 $\mathcal{E}$（或同态 $\sigma : W_F \to {}^{L}G$）的 Hecke 本征函数。

% [Tutor Note: 推导补足。由 $\mathbb{F}_q$-点上的本征性推广到 $m$ 次闭点，关键是这两步：(1) 把 $\alpha_V^{\boxtimes m}$ 对称化为 $\alpha_V^{(m)}$，使 Hecke 函子生活在 $\mathrm{Sym}^m C$ 上，于是算子落在「除子」$D(x)$ 而不是单个点上；(2) 上一条显示的迹等式——它是 (8.4) 中 $p_*$ 与 $(-)^{S_m}$ 的直接推论：$\mathrm{Fr}_x$ 在 $m$ 个共轭点 $y_1,\dots,y_m$ 上循环置换，对称化后正好把「$m$ 个点上各自的 $\mathrm{Fr}_{y_i}$ 之积」压缩成 $\mathrm{Sym}^m C$ 上单点 $D(x)$ 处的 $\mathrm{Fr}_{D(x)}$，而局部系统 $(V_{\mathcal{E}})^{\boxtimes m}$ 的 $S_m$-不变部分在该点的迹恰是 $\mathrm{Tr}(\mathrm{Fr}_x,V_{\mathcal{E}})$。没有这一步，就只能得到 (8.3) 的 $\mathbb{F}_q$-点版本。 来源：原书 p.99 上下文（A）。]

「朴素」的 Hecke 本征层性质 (8.1) 本身并不蕴含那些剩余域为 $\mathbb{F}_q$（我们曲线的定义域）的非平凡扩张的闭点处的 Hecke 本征函数性质，这一事实多少有些出人意料。然而，为保证 Hecke 本征函数性质处处成立所需的 $S_2$-等变性是一个非常自然的条件。事实上，它是 [FGV] 第 1.4 节所引入的如下 Hecke 函子的一般交换性条件的特殊情形。

对 $V, W\in \mathrm{Rep}({}^{L}G)$，设 $H_V$ 与 $H_W$ 是对应的 Hecke 函子，从 $\Bun_G$ 上层之导出范畴作用到 $C\times \Bun_G$ 上层之导出范畴。于是我们有迭代函子 $H_V\circ H_W$，从 $\Bun_G$ 上层之导出范畴作用到 $C\times C\times \Bun_G$ 上层之导出范畴。给定 Hecke 本征层 $(\mathcal{F}, (\alpha_V))$，我们有同构
\[
\alpha_V\circ \alpha_W : (H_V\circ H_W)(\mathcal{F}) \simeq \mathcal{E}_V\boxtimes \mathcal{E}_W\boxtimes \mathcal{F}.
\]
另一方面，在 $C\times C\setminus\Delta$ 上我们有自然的等同
\[
(H_V\circ H_W)(\mathcal{F})|_{C\times C\setminus\Delta} \simeq \sigma^*\circ (H_W\circ H_V)(\mathcal{F})|_{C\times C\setminus\Delta},
\]
其中 $\sigma$ 是 $C\times C\setminus\Delta$ 上的换位。交换性条件就是下图交换：
\[
\begin{array}{ccc}
(H_V\circ H_W)(\mathcal{F})|_{C\times C\setminus\Delta} & \xrightarrow{\ \alpha_V\circ\alpha_W\ } & \mathcal{E}_V\boxtimes \mathcal{E}_W\boxtimes \mathcal{F}|_{C\times C\setminus\Delta}\\[10pt]
\Big\downarrow & & \Big\downarrow\\[10pt]
\sigma^*\circ (H_W\circ H_V)(\mathcal{F})|_{C\times C\setminus\Delta} & \xrightarrow{\ \sigma^*(\alpha_W\circ\alpha_V)\ } & \sigma^*(\mathcal{E}_W\boxtimes \mathcal{E}_V)\boxtimes \mathcal{F}|_{C\times C\setminus\Delta}
\end{array}
\]
若 $V = W$，我们就得到上述 $S_2$-等变条件。

% [存疑: 原文第 98 页把该本征值层写作 $\mathcal{V}_{\mathcal{E}}$，而第 99 页起改写作 $\mathcal{E}_V$、$\mathcal{E}_W$（字母位置与字体均有变化）。两处记号不一致，译文各按原文保留，未作统一。 来源：原书 p.99 英文原文；存疑待核（A）。]

为解释这一交换性条件的含义，我们从 [MV] 回忆几何 Satake 等价：它是在支撑在固定点 $x\in C$ 的 Hecke 函子范畴（即仿射 Grassmann 簇上的等变反常层范畴）与范畴 $\mathrm{Rep}({}^{L}G)$ 之间的等价。这是张量范畴的等价，意思是：除与两个范畴中的张量积相容之外，它还与交换性约束和结合性约束相容。在前一个范畴上，交换性约束（按照 V. Drinfeld 的做法）被定义为定义在 $C$ 的不同点上的 Hecke 函子的换位在那些点重合时的某个极限。

% [Tutor Note: 跳步补足。几何 Satake 等价（Mirković--Vilonen [MV]）的内容是：$\mathrm{Gr}_G = G((t))/G[[t]]$ 上的 $G[[t]]$-等变反常层范畴，以卷积为张量积，等价于 $\mathrm{Rep}({}^{L}G)$ 的张量范畴；不可约等变反常层对应不可约表示（由 $G[[t]]$-轨道即 coweight 的支配性标记，这正是 Satake 同构的几何版本）。原文强调「与交换性约束相容」并非空话：张量范畴等价必须同时匹配结合约束与交换约束，而后者在 Grassmann 侧不是显然的，Drinfeld 用「两点上的 Hecke 函子在两点重合时的换位极限」来定义它。下文的 Tannakian 论证正是要说明：只给结合性（幺半结构）不足以得到群，必须再给交换约束。 来源：原书 p.99 上下文（A）。]

（正则）Hecke 本征层的概念可以看作「从支撑在一点 $x\in C$ 的 Hecke 函子范畴到 $\mathrm{Rep}({}^{L}G)$ 的纤维函子」这一概念的自然推广：我们允许点 $x$ 沿曲线移动。从这个观点看，要求同构 $\alpha_V$ 与张量积结构及结合性相容，相当于要求该纤维函子
% ==== 原文 PDF 第 100 页（印刷页 100）====
是幺半函子（monoidal functor，即与张量积及结合性约束相容的函子）。但我们从 Tannakian 理论知道，这还不足以建立一个张量范畴与某个代数群的表示范畴之间的等价。为此我们还需要纤维函子与交换性约束相容。由于支撑在一点的 Hecke 函子范畴上的交换性约束表现为支撑在 $C$ 的两个不同点上的两个 Hecke 函子的换位之极限，交换性约束本身就表现为上述交换性条件在两个点重合时的极限。

因此我们看到，要求 Hecke 本征层满足交换性条件是相当自然的。我们上面观察到的一个有趣事实是：这个条件的一部分（$V = W$ 的情形）对于保证「在正特征中工作时，与 Hecke 本征层相联系的函数是曲线所有闭点处的 Hecke 本征函数」也是必要的。

\subsection{回到 \texorpdfstring{$SL_2$}{SL2}}

正如上一节所见，双商 (8.2) 的几何对应物是 $C$ 上 $G$-丛的模叠 $\Bun_G$。事实上，(8.2) 就是 $\Bun_G$ 的 $\mathbb{F}_q$-点之集。因此，$\Bun_G$ 上的 Hecke 本征层给出 (8.2) 上的 Hecke 本征函数，如上所述。

另一方面，我们在 7.3 节看到，为恰当理解 $G = SL_2$ 的（非分歧）自守表示的 L-包，我们需要考虑更一般的双商 (7.12)，其中 $K_x = K'_x$ 或 $K''_x$，且对除有限多个闭点 $x\in C$ 之外的所有闭点有 $K_x = K'_x = SL_2[[t]]$。若所有 $K_x = SL_2[[t]]$，则 (7.12) 是 $\Bun_{SL_2}$ 的 $\mathbb{F}_q$-点之集。更一般的商 (7.12) 是什么？答案很清楚：它们是 $\Bun_{SL_2}$ 的「improper」版本的 $\mathbb{F}_q$-点之集；即行列式为线丛 $\mathcal{L} = \mathcal{O}(D)$ 的秩二向量丛的模叠 $\Bun^{\mathcal{L}}_{SL_2}$。这里 $D$ 是满足 $K_x = K''_x$ 的点之集，我们把它视为 $C$ 上的有效除子。

我们在曲线定义在 $\mathbb{C}$ 上的情形已经遇到过这些模叠，见 3.8、5.1.2 与 5.2.3 节。当时我们指出，若 $\mathcal{L} = \mathcal{L}'\otimes \mathcal{N}^2$，其中 $\mathcal{N}$ 是 $C$ 上的线丛，则可通过把秩二向量丛与 $\mathcal{N}$ 作张量积，把 $\Bun^{\mathcal{L}}_{SL_2}$ 与 $\Bun^{\mathcal{L}'}_{SL_2}$ 等同起来。因此 $\Bun^{\mathcal{L}}_{SL_2}$ 实际上不依赖于 $\mathcal{L}$，而依赖于它在 Picard 群 $\mathrm{Pic}(C)$（即 $C$ 的 Picard 概形 $\mathrm{Pic}$ 的 $\mathbb{C}$-点之集）对平方子群取商所得的商中的像。这个商同构于 $\mathbb{Z}_2$，因此在这一情形下 improper 分支 $\Bun^{\mathcal{L}}_{SL_2}$ 在非典范同构的意义下唯一（此时我们可以取 $\mathcal{L} = \mathcal{O}(p)$，其中 $p$ 是 $C$ 的一点）。

现在考虑定义在 $\mathbb{F}_q$ 上的曲线 $C$。这里 improper 模叠 $\Bun^{\mathcal{L}}_{SL_2}$ 同样在同构意义下由商 $\mathrm{Pic}(\mathbb{F}_q)/\mathrm{Pic}(\mathbb{F}_q)^2$ 分类。但现在这个商大得多。为更精确地描述它，我们回忆 [Se]：交换类域论把 $\mathrm{Pic}(\mathbb{F}_q)$ 与函数域 $F$（$C$ 的函数域）的 Weil 群 $W_F$ 的最大非分歧阿贝尔商等同起来。换言之，$\mathrm{Pic}(\mathbb{F}_q)$ 同构于 $F$ 的最大非分歧阿贝尔扩张 $F^{\mathrm{ab},\mathrm{un}}$ 的 Galois 群的稠密子群；该扩张的定义方式与 $F$ 的 Weil 群相同。也就是说，它是 $\mathbb{Z}\subset \hat{\mathbb{Z}}$ 在同态 $\mathrm{Gal}(F^{\mathrm{ab},\mathrm{un}}/F)\to \mathrm{Gal}(\overline{\mathbb{F}}_q/\mathbb{F}_q) = \hat{\mathbb{Z}}$ 下的原像。于是我们得到：$\mathrm{Pic}(\mathbb{F}_q)/\mathrm{Pic}(\mathbb{F}_q)^2$ 是 $\mathrm{Gal}(F^{\mathrm{ab},\mathrm{un}}/F)$ 的使所有
% ==== 原文 PDF 第 101 页（印刷页 101）====
元素均为 $2$ 阶的最大商。它也是 $F$ 的非分歧二次扩张之群的对偶群。事实上，每个这样的扩张 $E/F$ 都给出 $\mathrm{Pic}(\mathbb{F}_q)$ 的一个二次特征标，它通过 $\mathrm{Pic}(\mathbb{F}_q)/\mathrm{Pic}(\mathbb{F}_q)^2$ 分解。

% [Tutor Note: 推导补足。这一步用到的类域论对应是：$F$ 的有限非分歧阿贝尔扩张 $\leftrightarrow$ $\mathrm{Gal}(\overline{\mathbb{F}}_q/\mathbb{F}_q)=\hat{\mathbb{Z}}$ 的连续商 $\leftrightarrow$ $\hat{\mathbb{Z}}$ 的开子群，而 $\hat{\mathbb{Z}}$ 的开子群恰是 $n\hat{\mathbb{Z}}$，故非分歧阿贝尔扩张全体由「取 $n$ 次扩域」标记，其 Galois 群 $\cong \mathbb{Z}/n$。于是 $\mathrm{Pic}(\mathbb{F}_q)\cong W_F^{\mathrm{ab},\mathrm{un}}$ 的对偶（离散）刻画给出：$\mathrm{Pic}(\mathbb{F}_q)/\mathrm{Pic}(\mathbb{F}_q)^2$ 对偶于「非分歧二次扩张之集」，而后者正是 $\mathrm{Pic}(\mathbb{F}_q)$ 的 $2$-挠部分的对偶，故该商是 $\mathbb{F}_2$-向量空间。原文说它「是 $F$ 的非分歧二次扩张之群的对偶群」，严格说该「群」指 $\mathrm{Pic}(\mathbb{F}_q)/2\mathrm{Pic}(\mathbb{F}_q)$，其元素为二次扩张/二次特征标；在正特征 $p\neq 2$ 时这与 $\mathrm{Pic}(\mathbb{F}_q)/\mathrm{Pic}(\mathbb{F}_q)^2$ 一致，$p=2$ 时须留神 $2$-挠与平方子群不再同构。 来源：原书 p.101 上下文（A）。]

我们选取这个群的一个代表 $\mathcal{L} = \mathcal{O}(D)$，其中 $D$ 是 $C$ 的闭点之集的一个子集。于是我们有 $C$ 上行列式为 $\mathcal{L}$ 的秩二向量丛的代数模叠 $\Bun^{\mathcal{L}}_{SL_2}$，其 $\mathbb{F}_q$-点之集就是按上述子群 $K_x$ 的选取所得的双商 (7.12)。在几何理论中，这个集合上 Hecke 本征函数的概念变成（正则）Hecke 本征层的概念，其定义方式与（对应于 $\mathcal{L} = \mathcal{O}$ 的）proper 模叠 $\Bun_{SL_2}$ 的情形相同。

\subsection{从 \texorpdfstring{$\mathbb{C}$}{C} 上的曲线到 \texorpdfstring{$\mathbb{F}_q$}{Fq} 上的曲线}

回到 $\mathbb{C}$ 上的曲线 $C$，并选取一个 $PGL_2$-局部系统 $\mathcal{E}$，其结构群约化到 $O_2\subset PGL_2$，但不约化到它的真子群。由于 $\mathcal{E}$ 来自一个不可约秩二局部系统，我们预期 $\Bun_{SL_2}$ 上以 $\mathcal{E}$ 为本征值的正则 Hecke 本征层（6.5 节意义下）的范畴有一个不可约对象（在同构意义下）。设 $\mathcal{F}$ 是 $\Bun_{SL_2}$ 上相应的 $\CD$-模。在第 5 节我们讨论过与 $\mathcal{F}$ 相对应的 A-膜 $\mathcal{A}$，它由奇异 Hitchin 纤维上的一个秩一酉局部系统表示，而该纤维有两个不可约分支。我们观察到 $\mathcal{A}$ 分裂为两个 A-膜 $\mathcal{A}_+$ 与 $\mathcal{A}_-$，分别支撑在 Hitchin 纤维的两个不可约分支上。因此我们预期 $\CD$-模 $\mathcal{F}$ 也分裂为直和
\begin{equation}\tag{8.5}
\mathcal{F} = \mathcal{F}_+ \oplus \mathcal{F}_-,
\end{equation}
即 $\Bun_{SL_2}$ 上两个不可约 $\CD$-模的直和，对应于奇异 Hitchin 纤维上的两个 A-膜。此外，由于 A-膜 $\mathcal{A}_\pm$ 关于 't Hooft 算子是分数本征膜，我们预期层 $\mathcal{F}_\pm$ 满足 6.5 节所引入的分数 Hecke 性质。

这引导我们假定同样的现象对有限域 $\mathbb{F}_q$ 上的曲线也应发生。也就是说，对应于 $\mathbb{F}_q$ 上曲线 $C$ 的 $\ell$-进局部系统 $\mathcal{E}$（如上文所讨论的那一类）的正则 Hecke 本征层 $\mathcal{F}$，也应分裂为直和 (8.5)。此外，这些层应满足 6.5 节所引入的分数 Hecke 性质，从而给出分数 Hecke 本征层的范畴。其次，在有限域上曲线的设定中，我们可以从 $\Bun_{SL_2}$ 上的 $\ell$-进反常层过渡到函数。于是层 $\mathcal{F}_\pm$ 中的每一个都应给出双商 (8.2)（即 $\Bun_{SL_2}$ 的 $\mathbb{F}_q$-点之集）上的函数 $f_\pm$。层 $\mathcal{F}_\pm$ 的分数 Hecke 性质转化为相应函数 $f_\pm$ 的某个性质。

于是，我们从 A-膜出发，最终得到满足分数 Hecke 性质的自守函数。示意地说，这一过渡如下：
\[
\begin{array}{ccccccc}
\boxed{\text{A-branes}} & \xRightarrow{\ \text{over }\mathbb{C}} &
\boxed{\CD\text{-modules}} & \xRightarrow{\ \text{over }\mathbb{C}} &
\boxed{\begin{smallmatrix}\text{perverse}\\ \text{sheaves}\end{smallmatrix}} & \xRightarrow{\ \text{over }\mathbb{F}_q} &
\boxed{\text{functions}}
\end{array}
\]
下文将看到，分数 Hecke 性质特别意味着：不仅 $f_+ + f_-$ 是通常意义下的 Hecke 本征函数，而且 $f_+ - f_-$ 也是 Hecke
% ==== 原文 PDF 第 102 页（印刷页 102）====
本征函数，只不过它是关于另一个同态 $\sigma' : W_F \to PGL_2$ 的本征函数。我们将证明 $\sigma'$ 确实存在，并且事实上典范地附属于原来的同态 $\sigma$。这将为我们的预言提供第一次一致性检验。

\subsection{分数 Hecke 性质}

设 $C$ 是 $\mathbb{F}_q$ 上的曲线，$\mathcal{E}$ 是 $C$ 上的内窥 $\ell$-进 $PGL_2$-局部系统（对应于非分歧同态 $\sigma : W_F \to PGL_2$）。这意味着它的结构群约化到 $O_2$，但不约化到一个真子群。于是 $\mathcal{E}$ 的自同构群（等价地，$\sigma$ 的像的中心化子）是 $\mathbb{Z}_2$。设 $D$ 是 $C$ 的闭点的有限集。用 $\mathcal{F}^D$ 表示 $\Bun^{\mathcal{O}(D)}_{SL_2}$ 上以 $\mathcal{E}$ 为「本征值」的正则 Hecke 本征层（6.5 节意义下）。受我们在 $\mathbb{C}$ 上曲线类似情形中关于 A-膜的结果的启发，我们猜想 $\mathcal{F}^D$ 分裂为直和
\begin{equation}\tag{8.6}
\mathcal{F}^D = \mathcal{F}^D_+ \oplus \mathcal{F}^D_-,
\end{equation}
即反常层 $\mathcal{F}^D_\pm$ 的直和，它们关于 $\mathcal{E}$ 满足 6.5 节所引入的如下分数 Hecke 性质（因此我们也称它们为分数 Hecke 本征层）：
\begin{equation}\tag{8.7}
\alpha_+ : H_W(\mathcal{F}^D_+) \xrightarrow{\ \sim\ } \bigl(\det U_{\mathcal{E}}\boxtimes \mathcal{F}^D_+\bigr)\oplus \bigl(U_{\mathcal{E}}\boxtimes \mathcal{F}^D_-\bigr),
\end{equation}
\begin{equation}\tag{8.8}
\alpha_- : H_W(\mathcal{F}^D_-) \xrightarrow{\ \sim\ } \bigl(U_{\mathcal{E}}\boxtimes \mathcal{F}^D_+\bigr)\oplus \bigl(\det U_{\mathcal{E}}\boxtimes \mathcal{F}^D_-\bigr).
\end{equation}
这里 $W$ 是 $PGL_2$ 的伴随表示，我们对 $C$ 上秩三局部系统 $W_{\mathcal{E}}$ 按其上自同构群 $\mathbb{Z}_2$ 的作用作如公式 (6.8) 那样的分解：
\begin{equation}\tag{8.9}
W_{\mathcal{E}} = \bigl(\det U_{\mathcal{E}}\otimes I\bigr)\oplus \bigl(U_{\mathcal{E}}\otimes S\bigr),
\end{equation}
其中 $I$ 与 $S$ 分别是 $\mathbb{Z}_2$ 的平凡表示与符号表示，而 $\det U_{\mathcal{E}}$ 与 $U_{\mathcal{E}}$ 是 $C$ 上的秩一与秩二局部系统，定义如下。回忆按我们的假设，$PGL_2$-局部系统 $\mathcal{E}$ 约化到 $O_2$，所以我们把它视为 $O_2$-局部系统。然后我们置
\[
U_{\mathcal{E}} = \mathcal{E}\times_{O_2} U,
\]
其中 $U$ 如前所述是 $O_2$ 的定义二维表示。

% [Tutor Note: 推导补足。(8.7)--(8.9) 合起来说明「分数」二字何来。因为 $\mathcal{E}$ 的自同构群是 $\mathbb{Z}_2$，$W_{\mathcal{E}}$ 才按 $I$ 与 $S$ 两个同型分量分解：$I$-分量（对应本征值 $a_x$，即 $\det U_{\mathcal{E}}$ 的迹）把 $\mathcal{F}^D_+$ 送到 $\mathcal{F}^D_+$、把 $\mathcal{F}^D_-$ 送到 $\mathcal{F}^D_-$；而 $S$-分量（符号表示，对应本征值 $b_x$）把 $\mathcal{F}^D_+$ 与 $\mathcal{F}^D_-$ 互换。因此单个 $\mathcal{F}^D_\pm$ 不是 $H_W$ 的本征对象（本征对象须是「本征值乘自身」），$H_W$ 只能把它送到两者的直和——这正是 (8.7)(8.8) 右端的结构，也正是「分数」本征性的含义。物理上：'t Hooft 算子在奇异 Hitchin 纤维的两个分支之间交换，单个分支上的 A-膜只能在两个分支的张成空间里是本征的。 来源：原书 p.102 上下文（A）。]

\subsection{分数 Hecke 本征函数}

现在我们分析公式 (8.7) 与 (8.8) 对与 $\mathcal{F}^D_\pm$ 相联系的函数的推论，
\[
f^D_\pm = \mathrm{f}_{\mathcal{F}^D_\pm,\mathbb{F}_q},
\]
这些函数定义在集合 $\Bun^{\mathcal{O}(D)}_{SL_2}(\mathbb{F}_q)$ 上，而该集合同构于双商 (7.12)。公式 (8.6) 蕴含
\begin{equation}\tag{8.10}
f^D = f^D_+ + f^D_-,
\end{equation}
其中 $f^D = \mathrm{f}_{\mathcal{F}^D,\mathbb{F}_q}$ 是与正则 Hecke 本征层 $\mathcal{F}^D$ 相联系的 $\Bun^{\mathcal{O}(D)}_{SL_2}(\mathbb{F}_q)$ 上的函数。

为简化记号，下文中在没有歧义时我们将省略上标 $D$。
% ==== 原文 PDF 第 103 页（印刷页 103）====
把 Hecke 对应限制到 $x\times \Bun^{\mathcal{O}(D)}_{SL_2}$（其中 $x\in C(\mathbb{F}_q)$）并在那里的 $\mathbb{F}_q$-点上计算 Frobenius 的迹，我们从公式 (8.7) 与 (8.8) 得到：函数 $f_\pm$ 满足如下性质
\begin{equation}\tag{8.11}
T_{W,x}\cdot \begin{pmatrix} f_+\\ f_-\end{pmatrix} = \begin{pmatrix} a_x & b_x\\ b_x & a_x\end{pmatrix}\begin{pmatrix} f_+\\ f_-\end{pmatrix}, \qquad x\in C(\mathbb{F}_q),
\end{equation}
其中
\[
a_x = \mathrm{Tr}(\mathrm{Fr}_x, \det U_{\mathcal{E}}) = \det(\sigma_x(\mathrm{Fr}_x), U),
\]
\[
b_x = \mathrm{Tr}(\mathrm{Fr}_x, U_{\mathcal{E}}) = \mathrm{Tr}(\sigma_x(\mathrm{Fr}_x), U).
\]
这里我们把 $\sigma_x$ 视为同态 $W_{F_x}\to O_2$。为计算这些数，我们回忆 7.2 节对 Frobenius 共轭类的描述。我们发现，若 $x$ 分裂，则 $\sigma_x(\mathrm{Fr}_x)$ 在 $O_2$ 中的共轭类包含矩阵
\[
\begin{pmatrix}
\dfrac{\mu(\mathrm{Fr}_{y_1(x)})}{\mu(\mathrm{Fr}_{y_2(x)})} & 0\\[10pt]
0 & \dfrac{\mu(\mathrm{Fr}_{y_2(x)})}{\mu(\mathrm{Fr}_{y_1(x)})}
\end{pmatrix},
\]
而若 $x$ 非分裂，则包含矩阵
\[
\begin{pmatrix} 1 & 0\\ 0 & -1\end{pmatrix}.
\]
因此我们发现
\begin{equation}\tag{8.12}
a_x = \begin{cases} 1, & \text{if } x \text{ is split},\\ -1, & \text{if } x \text{ is non-split},\end{cases}
\end{equation}
\begin{equation}\tag{8.13}
b_x = \begin{cases} \dfrac{\mu(\mathrm{Fr}_{y_1(x)})}{\mu(\mathrm{Fr}_{y_2(x)})} + \dfrac{\mu(\mathrm{Fr}_{y_2(x)})}{\mu(\mathrm{Fr}_{y_1(x)})}, & \text{if } x \text{ is split},\\[10pt] 0, & \text{if } x \text{ is non-split}.\end{cases}
\end{equation}

% [Tutor Note: 易错预警。(8.12) 中 $a_x = \pm 1$ 的来源是 $\det$ 而非 $\mathrm{Tr}$：分裂情形对角元之积 $\frac{\mu(\mathrm{Fr}_{y_1})}{\mu(\mathrm{Fr}_{y_2})}\cdot\frac{\mu(\mathrm{Fr}_{y_2})}{\mu(\mathrm{Fr}_{y_1})}=1$，非分裂情形 $\det\mathrm{diag}(1,-1)=-1$；所以 $a_x$ 只记录「分裂/非分裂」这一 $\mathbb{Z}_2$-信息，它正是 (8.9) 中 $I$-分量（$\mathbb{Z}_2$ 的平凡表示）的贡献。而 (8.13) 的 $\mathrm{Tr}$ 恒等于对角元之和，对应 $S$-分量（符号表示）的贡献，故分裂时非零、非分裂时为零（非分裂时 $\mathrm{diag}(1,-1)$ 的迹为 $0$）。注意原文这里把 $\mu(\mathrm{Fr}_{y_1(x)})$ 写在分子、$\mu(\mathrm{Fr}_{y_2(x)})$ 写在分母，与 §7.2 中 $\mu$ 是 $W_{F_x}$ 的一维特征标、$y_1,y_2$ 是 $x$ 在二次扩张中的两个原像的约定一致，抄写时切勿上下颠倒。 来源：原书 p.103 上下文（A）。]

公式 (8.11) 蕴含和 $f_+ + f_-$ 是 Hecke 算子的本征函数
\begin{equation}\tag{8.14}
T_{W,x}\cdot (f_+ + f_-) = (a_x + b_x)(f_+ + f_-), \qquad x\in C(\mathbb{F}_q),
\end{equation}
其中
\begin{equation}\tag{8.15}
\begin{aligned}
a_x + b_x &= \mathrm{Tr}(\mathrm{Fr}_x, W_{\mathcal{E}}) = \mathrm{Tr}(\sigma_x(\mathrm{Fr}_x), W)\\
&= \begin{cases} 1 + \dfrac{\mu(\mathrm{Fr}_{y_1(x)})}{\mu(\mathrm{Fr}_{y_2(x)})} + \dfrac{\mu(\mathrm{Fr}_{y_2(x)})}{\mu(\mathrm{Fr}_{y_1(x)})}, & \text{if } x \text{ is split},\\[10pt] -1, & \text{if } x \text{ is non-split}.\end{cases}
\end{aligned}
\end{equation}
于是，公式 (8.14) 表达了与层 $\mathcal{F} = \mathcal{F}_+ \oplus \mathcal{F}_-$ 相联系的函数 $f = f_+ + f_-$ 关于 $\mathcal{E}$（或 $\sigma$）的通常 Hecke 性质。

但我们还发现，差 $f_+ - f_-$ 是具有另一组本征值的 Hecke 本征函数；即
\begin{equation}\tag{8.16}
T_{W,x}\cdot (f_+ - f_-) = (a_x - b_x)(f_+ - f_-), \qquad x\in C(\mathbb{F}_q),
\end{equation}
% ==== 原文 PDF 第 104 页（印刷页 104）====
其中
\begin{equation}\tag{8.17}
a_x - b_x = \begin{cases} 1 - \dfrac{\mu(\mathrm{Fr}_{y_1(x)})}{\mu(\mathrm{Fr}_{y_2(x)})} - \dfrac{\mu(\mathrm{Fr}_{y_2(x)})}{\mu(\mathrm{Fr}_{y_1(x)})}, & \text{if } x \text{ is split},\\[10pt] -1, & \text{if } x \text{ is non-split},\end{cases}
\end{equation}
然而，我们必须记住 Hecke 性质 (8.16) 只对 $C$ 的 $\mathbb{F}_q$-点成立。事实上，我们是从几何 Hecke 性质 (8.7)--(8.8) 出发的。Hecke 对应联系 $\Bun^{\mathcal{O}(D)}_{SL_2}$ 与 $C\times \Bun^{\mathcal{O}(D)}_{SL_2}$。函数 $f_\pm$ 由 $\mathcal{F}_\pm$ 在 $\Bun^{\mathcal{O}(D)}_{SL_2}$ 的 $\mathbb{F}_q$-点上取 Frobenius 的迹得到。为得到它们的 Hecke 性质，我们需要考虑这两个模叠的 $\mathbb{F}_q$-点之集上的 Hecke 对应。这样能达到的 Hecke 算子只是对应于 $C$ 的 $\mathbb{F}_q$-点的那些。Hecke 算子由此产生的作用由方程 (8.11) 表达。

这个公式并不唯一确定函数 $f_+ - f_-$。只有当它是 $C$ 的所有闭点（而不仅是其 $\mathbb{F}_q$-点）上的 Hecke 本征函数时，它才会被唯一确定（至少对于我们正在考虑的这类一般的 $\sigma$，即对应于 $W_F$ 的不可约二维表示的那些）。

为纳入剩余域为 $\mathbb{F}_{q^m}$（$m > 1$）的 $C$ 的闭点，我们需要考虑 $C$ 的第 $m$ 个对称幂上的更一般的 Hecke 对应 $H_W^{(m)}$（并去掉成对对角线之并 $\Delta$）。这要求同构 $\alpha_\pm$ 满足附加的 $S_2$-等变条件，类似于正则 Hecke 本征层情形中的条件。这个条件的定义方式与 8.2 节完全相同。

设这一 $S_2$-等变条件成立，我们得到同构
\[
\begin{array}{r@{\ }l}
\alpha_+^{(2)} : H_W^{(2)}(\mathcal{F}_+) \xrightarrow{\ \sim\ } & \bigl((\det U_{\mathcal{E}})^{(2)}\oplus U_{\mathcal{E}}^{(2)}\bigr)\big|_{\mathrm{Sym}^2 C\setminus\Delta}\boxtimes \mathcal{F}_+\\[4pt]
& \oplus\ \mathrm{Sym}\bigl(\det U_{\mathcal{E}}\boxtimes U_{\mathcal{E}}\oplus U_{\mathcal{E}}\boxtimes \det U_{\mathcal{E}}\bigr)\big|_{\mathrm{Sym}^2 C\setminus\Delta}\boxtimes \mathcal{F}_-,
\end{array}
\]
\[
\begin{array}{r@{\ }l}
\alpha_-^{(2)} : H_W^{(2)}(\mathcal{F}_-) \xrightarrow{\ \sim\ } & \mathrm{Sym}\bigl(\det U_{\mathcal{E}}\boxtimes U_{\mathcal{E}}\oplus U_{\mathcal{E}}\boxtimes \det U_{\mathcal{E}}\bigr)\big|_{\mathrm{Sym}^2 C\setminus\Delta}\boxtimes \mathcal{F}_+\\[4pt]
& \oplus\ \bigl((\det U_{\mathcal{E}})^{(2)}\oplus U_{\mathcal{E}}^{(2)}\bigr)\big|_{\mathrm{Sym}^2 C\setminus\Delta}\boxtimes \mathcal{F}_-.
\end{array}
\]
这里我们使用了记号 (8.4)。

设这个条件满足。我们从 8.2 节回忆：$C$ 的任一满足 $\deg(x)=m$ 的闭点 $x$（$\deg(x)$ 的定义见 7.1 节）给出 $\mathrm{Sym}^m C$ 中的一个 $\mathbb{F}_q$-点 $D(x)$（一个 $m$ 次有效除子），并且对 $C$ 上任一局部系统 $\mathcal{E}$ 有
\[
\mathrm{Tr}\bigl(\mathrm{Fr}_{D(x)}, \mathcal{E}^{(m)}\bigr) = \mathrm{Tr}(\mathrm{Fr}_x, \mathcal{E}).
\]
特别地，设 $x$ 是 $C$ 的二次闭点，即其剩余域同构于 $\mathbb{F}_{q^2}$，并设 $D(x)$ 是 $\mathrm{Sym}^2 C$ 中相应的 $\mathbb{F}_q$-点。则我们有
\[
\mathrm{Tr}\bigl(\mathrm{Fr}_{D(x)}, \det U_{\mathcal{E}}^{(2)}\oplus U_{\mathcal{E}}^{(2)}\bigr) = \mathrm{Tr}\bigl(\mathrm{Fr}_{D(x)}, W_{\mathcal{E}}^{(2)}\bigr) = \mathrm{Tr}(\mathrm{Fr}_x, W_{\mathcal{E}}) = a_x + b_x,
\]

% [存疑: 本式中原文写作 $\det U_{\mathcal{E}}^{(2)}$（$\det$ 与上标 $(2)$ 之间没有括号），而前面 $\alpha_\pm^{(2)}$ 的两个同构式中同一对象写作 $(\det U_{\mathcal{E}})^{(2)}$。按记号 (8.4)，$(\det U_{\mathcal{E}})^{(2)}$ 才是对称化后所得的层；两处译文均照原文各自的写法保留，未作统一。 来源：原书 p.104 英文原文；存疑待核（A）。]
% ==== 原文 PDF 第 105 页（印刷页 105）====
但
\[
\mathrm{Tr}\bigl(\mathrm{Fr}_{D(x)}, \mathrm{Sym}\bigl(\det U_{\mathcal{E}}\boxtimes U_{\mathcal{E}}\oplus U_{\mathcal{E}}\boxtimes \det U_{\mathcal{E}}\bigr)\bigr) = 0.
\]
因此同构 $\alpha_\pm^{(2)}$ 蕴含 $f_+$ 与 $f_-$ 都是（通常意义下的）以 $(a_x + b_x)$ 为本征值的 Hecke 本征函数，于是我们有
\begin{equation}\tag{8.18}
T_{W,x}\cdot (f_+ - f_-) = (a_x + b_x)(f_+ - f_-), \qquad x\in C,\ \deg(x) = 2.
\end{equation}

% [Tutor Note: 推导补足。此处被跳过的一步是 $\mathbb{Z}_2$-同型分解在 $S_2$-对称化下的表现：由 (8.9) 知 $W_{\mathcal{E}} = \mathrm{det}U_{\mathcal{E}}\oplus (U_{\mathcal{E}}\otimes S)$，对 $\mathrm{Sym}^2$ 有 $\mathrm{Sym}^2 W_{\mathcal{E}} = \mathrm{Sym}^2(\mathrm{det}U_{\mathcal{E}})\oplus (\mathrm{det}U_{\mathcal{E}}\otimes U_{\mathcal{E}}\otimes S)\oplus \mathrm{Sym}^2(U_{\mathcal{E}}\otimes S)$；$\mathbb{Z}_2$-不变部分正是 $(\det U_{\mathcal{E}})^{(2)}\oplus U_{\mathcal{E}}^{(2)}$（前一项来自 $I$-$I$ 的对称积，后一项来自 $S$-$S$），而交叉项 $\det U_{\mathcal{E}}\otimes U_{\mathcal{E}}\otimes S$ 带符号表示、在 $\mathbb{Z}_2$-不变部分外，故它所对应的项在取迹时贡献 $0$——这就是第二个显示为何等于零、以及为何 $f_+ - f_-$ 在二次闭点处反而满足 $(a_x+b_x)$ 而不是 $(a_x-b_x)$ 的原因。换句话说：二次点的 Frobenius 交换 $y_1,y_2$，恰好抹去了 (8.16) 中 $b_x$ 前的负号。 来源：原书 p.105 上下文（A）。]

下一步，我们以同样的方式分析 $C$ 的任意次 $m$ 的闭点处的行为。$S_2$-不变条件蕴含 $S_m$-不变条件，如 8.2 节所述。我们于是发现：$m$ 为奇数时函数 $f_+ - f_-$ 满足公式 (8.16)，而 $m$ 为偶数时它满足公式 (8.18)。因此我们有
\begin{equation}\tag{8.19}
T_{W,x}\cdot (f_+ - f_-) = \bigl(a_x + (-1)^m b_x\bigr)(f_+ - f_-), \qquad x\in C,\ \deg(x) = m.
\end{equation}
根据 $GL_2$ 的 Langlands 对应 [Dr1, Dr2]，这样的公式只有在如下情形才可能成立：$T_{W,x}$ 的本征值，即数 $a_x + (-1)^m b_x$，等于某个同态 $\sigma' : W_F \to PGL_2$ 的 $\mathrm{Tr}(\sigma'_x(\mathrm{Fr}_x, W)$。这给了我们一个机会来检验我们关于分解 (8.6) 存在的预言。

% [原文疑误: 原书 p.105 此处作 "equal to $\mathrm{Tr}(\sigma'_x(\mathrm{Fr}_x, W)$"，括号不配对：$\mathrm{Fr}_x$ 之后缺少右括号。译文照抄原样的记号，按 $\mathrm{Tr}(\sigma'_x(\mathrm{Fr}_x), W)$ 理解。 来源：原书 p.105 英文原文与译本对照（A）。]

事实上，构造具有这一性质的同态 $\sigma'$ 很容易。注意任一同态 $\sigma : W_F \to O_2$ 都可以用二次特征标 $\rho : W_F \to \mathbb{Z}_2 = \{\pm 1\}$ 来扭转，其中群 $\mathbb{Z}_2$ 与 $O_2$ 的中心等同。我们用 $\sigma \mapsto \sigma\otimes\rho$ 表示这一运算。

特别地，标量的二次扩张 $\mathbb{F}_{q^2}/\mathbb{F}_q$ 定义了二次扩张 $\mathbb{F}_{q^2}(C)/\mathbb{F}_q(C)$（回忆 $F = \mathbb{F}_q(C)$），从而定义了 $W_F$ 的一个二次特征标，我们将它记为 $\nu$。这个特征标由如下性质确定：
\begin{equation}\tag{8.20}
\nu(\mathrm{Fr}_x) = (-1)^m, \qquad x\in C,\ \deg(x) = m.
\end{equation}
令 $\sigma' = \sigma\otimes \nu$。则我们发现
\[
\mathrm{Tr}\bigl(\sigma'_x(\mathrm{Fr}_x), W\bigr) = \begin{cases} 1 + (-1)^m\dfrac{\mu(\mathrm{Fr}_{y_1(x)})}{\mu(\mathrm{Fr}_{y_2(x)})} + (-1)^m\dfrac{\mu(\mathrm{Fr}_{y_2(x)})}{\mu(\mathrm{Fr}_{y_1(x)})}, & \text{if } x \text{ is split},\\[10pt] -1, & \text{if } x \text{ is non-split},\end{cases}
\]
对所有的 $x\in C(\mathbb{F}_q)$ 成立。因此
\[
\mathrm{Tr}\bigl(\sigma'_x(\mathrm{Fr}_x), W\bigr) = a_x + (-1)^m b_x.
\]
于是我们得到：以 $a_x + (-1)^m b_x$ 为本征值（如公式 (8.19) 中那样）的 Hecke 本征函数确实存在！我们用 $f'$ 表示这个函数。

利用这个函数，我们现在可以解出 $f_+$ 与 $f_-$：
\begin{equation}\tag{8.21}
f_\pm = \frac{1}{2}(f\pm f').
\end{equation}
这些就是与我们猜想其存在的分数 Hecke 本征层 $\mathcal{F}_\pm$ 相对应的函数。\footnote{本文在 arXiv 上出现之后，我们从 S. Lysenko 处得知，层 $\mathcal{F}_\pm$ 可以利用他在 theta 提升方面的工作 [Ly1, Ly2] 构造出来。}

% [Tutor Note: 推导补足。(8.20)--(8.21) 是本节的核心构造，其内容是一个关于 $\mathbb{Z}_2$ 的 Fourier 变换：$f = f_+ + f_-$ 与 $f' = f_+ - f_-$ 是同一对函数 $f_\pm$ 的「和」与「差」，而 $(f,f')$ 分别是对 $\sigma$ 与 $\sigma' = \sigma\otimes\nu$ 的通常 Hecke 本征函数。于是分数 Hecke 本征函数 $f_\pm = \frac{1}{2}(f\pm f')$ 恰是「表示论意义下不可分解、但函数论意义下可分解」的对象。注意 $\sigma$ 与 $\sigma'$ 的自同构群都是 $\mathbb{Z}_2$、且附属于同一个二次扩张 $E/F$（见 8.8 节），故 $f,f'$ 落在同一个全局 L-包的两个成员所对应的空间里；这正是经典内窥中「L-包内两个表示由二次特征标 $\nu$ 区分」的几何对应物。$\nu$ 是 $\mathbb{F}_{q^2}/\mathbb{F}_q$ 诱导的，故 $\nu(\mathrm{Fr}_x) = (-1)^{\deg(x)}$ 只依赖于次数——这一「只依赖次数的符号」正是 (8.19) 中 $(-1)^m$ 的来源。 来源：原书 p.105 上下文（A）。]

% ==== 原文 PDF 第 106 页（印刷页 106）====
自然可以问：函数 $f_\pm$ 以及它们所满足的方程，其表示论意义是什么？这些方程由公式 (8.11)（对奇数次 $m$ 的 $x\in C$）给出，以及
\begin{equation}\tag{8.22}
T_{W,x}\cdot f_\pm = (a_x + b_x)f_\pm,
\end{equation}
（对偶数次 $m$ 的 $x\in C$）。

回忆 Hecke 本征函数，例如 $f = f_+ + f_-$ 与 $f' = f_+ - f_-$，可以解释为 $SL_2(\mathbb{A}_F)$ 的自守表示的矩阵系数，而它们的（正则）Hecke 性质是 Satake 同构的结果——该同构把球 Hecke 代数与 Langlands 对偶群 $PGL_2$ 的表示环等同起来。为分数 Hecke 本征函数 $f_\pm$ 以及方程 (8.11)、(8.22) 找到类似的解释将是有意义的。

% [原文疑误: 原书 p.106 此处将 "eigenfunctions" 误排为 "eienfunctions"（"the Hecke eienfunctions, such as $f = f_+ + f_-$"）；译文按正确拼写译为「Hecke 本征函数」。 来源：原书 p.106 英文原文与译本对照（A）。]

\subsection{improper Hecke 函子}

除了作用在 $\Bun^{\mathcal{O}(D)}_{SL_2}$ 上 $\CD$-模范畴上的「正常」Hecke 函子 $H_W$ 之外，还有「improper」Hecke 函子 $\tilde{H}_x$，它从 $\Bun^{\mathcal{O}(D)}_{SL_2}$ 上 $\CD$-模范畴作用到 $\Bun^{\mathcal{O}(D+x)}_{SL_2}$ 上 $\CD$-模范畴。它们通过两个模叠之间的 Hecke 对应来定义，该对应由秩二丛的对 $(\mathcal{M}, \mathcal{M}')$ 组成，其中 $\mathcal{M}\in \Bun^{\mathcal{O}(D)}_{SL_2}$，$\mathcal{M}'\in \Bun^{\mathcal{O}(D+x)}_{SL_2}$，且 $\mathcal{M}\subset \mathcal{M}'$ 作为凝聚层。这些函子是 7.4 节所引入的 improper Hecke 算子 $\tilde{T}_x$ 的范畴类比物。它们在 A-膜上对应的算子是 5.3 节所讨论的 improper 't Hooft 算子。

在公式 (5.38) 中我们计算了 improper 't Hooft 算子在膜 $\mathcal{A}_\pm$ 上的作用。基于这个公式，我们猜想 improper Hecke 算子应按如下方式作用在分数 Hecke 本征层 $\mathcal{F}^D_\pm$ 上：
\[
\tilde{H}_x(\mathcal{F}^D_+) \simeq \mathcal{F}^{D+x}_+ \oplus \mathcal{F}^{D+x}_-,\qquad
\tilde{H}_x(\mathcal{F}^D_-) \simeq \mathcal{F}^{D+x}_+ \oplus \mathcal{F}^{D+x}_-.
\]
若 $C$ 定义在 $\mathbb{C}$ 上，这应对所有点 $x\in C$ 成立；若 $C$ 定义在 $\mathbb{F}_q$ 上，则应对所有分裂点成立。

% [原文疑误: 原书 p.106 此处作 "Based in this formula"，介词应为 "on"；译文按 "Based on" 译出，不改变数学内容。 来源：原书 p.106 英文原文与译本对照（A）。]

这个公式表明，improper Hecke 函子未能建立内窥局部系统的情形下 $\Bun^{\mathcal{O}(D)}_{SL_2}$ 与 $\Bun^{\mathcal{O}(D+x)}_{SL_2}$ 上分数 Hecke 本征层范畴之间的等价。这可以看作 7.4 节所讨论的「作用在函数上的 improper Hecke 算子消失」这一现象的几何对应物；而正如我们已看到的，后者与内窥同态 $\sigma : W_F \to PGL_2$ 所对应的自守表示的全局 L-包的结构密切相关。

在正则 Hecke 本征层范畴的框架中，要看出 L-包现象的类比物会更为困难。事实上，对正则 Hecke 本征层 $\mathcal{F}^D = \mathcal{F}^D_+\oplus \mathcal{F}^D_-$，我们有
\begin{equation}\tag{8.23}
\tilde{H}_x(\mathcal{F}^D) \simeq V\otimes \mathcal{F}^{D+x},
\end{equation}
其中 $V$ 是二维向量空间。在 7.4 节所解释的经典设定中，$V$ 被迹 $\mathrm{Tr}(\tilde{\sigma}_x(\mathrm{Fr}_x), V)$ 取代。它恰在 $\sigma_x(\mathrm{Fr}_x)$
% ==== 原文 PDF 第 107 页（印刷页 107）====
是 (7.5) 的共轭类时消失，而这一消失体现了此情形下 L-包的非平凡结构。相反，在几何设定中，$V$ 本身作为乘子在公式 (8.23) 中出现。在此语境中什么应取代迹的消失，目前尚不清楚。

% [存疑: 上一句前半用 $\tilde{\sigma}_x(\mathrm{Fr}_x)$（$GL_2$ 的提升），后半却用不带 tilde 的 $\sigma_x(\mathrm{Fr}_x)$；按 8.8 节的约定，判别条件应针对 $\sigma_x$（$PGL_2$ 值）的共轭类，故此处原文可能漏写 tilde。译文照原文保留两种记号。 来源：原书 p.107 英文原文；存疑待核（A）。]

然而，在分数 Hecke 本征层范畴的框架中，图景更为透明。此时我们有一个含两个不可约对象 $\mathcal{F}^D_+$ 与 $\mathcal{F}^D_-$ 的范畴。函子 $\tilde{H}_x$ 把它们都送到 $\mathcal{F}^{D+x} = \mathcal{F}^{D+x}_+\oplus \mathcal{F}^{D+x}_-$，因此并没有在对应于 $D$ 与 $D+x$ 的范畴之间建立等价。

\subsection{与 \texorpdfstring{$\sigma$}{σ} 和 \texorpdfstring{$\sigma'$}{σ′} 相伴的 L-包}

上述讨论表明，对应于 $\sigma : W_F \to O_2\subset PGL_2$ 与 $\sigma' = \sigma\otimes\nu$（其中 $\nu$ 由公式 (8.20) 给出）的 $SL_2(\mathbb{A})$ 的表示 $\pi = \bigotimes'_x \pi_x$ 是相互联系的。比较对应于 $\sigma$ 与 $\sigma'$ 的 L-包（等价地，表示 $\pi$ 的重数）。

回忆对每个有限子集 $D\subset C$，我们有 $\pi$ 中 $\prod_x K_x$-不变向量构成的空间，其中 $x\in D$ 时 $K_x = K''_x$，否则 $K_x = K'_x$。若 $\pi$ 是自守的，则这个不变空间实现为双商 (7.12)（即 $\Bun^{\mathcal{O}(D)}_{SL_2}(\mathbb{F}_q)$）上 Hecke 本征函数空间。正如 7.3 节所解释的，这个空间或是一维或是零维。在前一情形，存在唯一的 Hecke 本征函数（相差一个标量），我们把它记为 $f^D_\sigma$。在后一情形，我们构造出的任何 Hecke 本征函数都必须为零。关于它是一维还是零维的判据由定理 7.1 给出（该定理遵循 [LL]），它相当于对全局 L-包的描述。

% [背景补充: 定理 7.1 即 Labesse--Langlands 对 $SL_2$ 的全局 L-包描述：一个局部几乎处处非分歧的 $SL_2$ 自守表示由 $PGL_2$-参数 $\sigma$ 与各处的局部分量共同决定，而「哪些局部 L-包成员能拼成一个全局自守表示」由局部不变量（这里是 $D$ 中各点关于二次扩张 $E/F$ 的分裂性）的奇偶性控制。原文给出的判据「$\#S$ 为偶则维数 $1$，为奇则维数 $0$」正是这一现象的精确陈述：$S$ 中每增加一个非分裂点，就多出一个「必须被抵消」的符号。 译者按（背景补充，无外部依据）。]

这个判据如下。考虑我们能将 $\sigma$ 提升为 $\tilde{\sigma} : W_F \to GL_2$ 并把 $\tilde{\sigma}$ 表示为 $\mathrm{Ind}_{W_E}^{W_F}\mu$ 的情形，其中 $E$ 是双重覆盖 $C'\to C$ 上的函数域，$\mu$ 是 $W_E$ 的特征标。则令 $S\subset D$ 为 $D$ 中在 $E$ 中非分裂的点之集。此时 $\Bun^{\mathcal{O}(D)}_{SL_2}(\mathbb{F}_q)$ 上关于 $\sigma$ 的 Hecke 本征函数空间的维数等于 $1$（若 $\#S$ 为偶数）、$0$（若 $\#S$ 为奇数）。

注意扩张 $E/F$ 对应于 $W_F$ 的如下二次特征标：它经由复合 $\kappa : W_F \xrightarrow{\ \sigma\ } O_2 \to \mathbb{Z}_2$ 得到。我们把 $E$ 称为与 $\sigma$ 相伴的域。

\begin{namedlemma}{8.1}
存在二次特征标 $\phi : W_E \to \mathbb{Z}_2$（其中 $E$ 与 $\sigma$ 相伴），使得 $\sigma' = \sigma\otimes\nu$ 等价于 $W_F$ 的二维表示 $\mathrm{Ind}_{W_E}^{W_F}(\mu\otimes\phi)$ 的射影化。
\end{namedlemma}

\textit{证明.} 选取 $\tilde{\tau}\in W_F$，它投影到 $W_F/W_E = \mathbb{Z}_2$ 的非平凡元素 $\tau$ 上。由 Cebotarev 定理，引理中所述条件等价于说 $\phi$ 满足
\[
\frac{\phi(\tilde{\tau}h\tilde{\tau}^{-1})}{\phi(h)} = \nu(h), \qquad h\in W_E.
\]
这样的 $\phi$ 的存在性可由 Hasse--Minkowski 定理推出，这一点由 B. Poonen 向我们解释。我们在此不给出他的论证，因为那会使我们离题太远。\hfill$\square$

% [Tutor Note: 推导补足。这个引理是 8.8 节的技术核心：$\sigma$ 与 $\sigma'$ 作为 $PGL_2$-参数只差一个二次特征标 $\nu$，问题在于 $\sigma'$ 是否仍来自某个 $GL_2$-提升，从而可用定理 7.1 的判据。条件 $\phi(\tilde{\tau}h\tilde{\tau}^{-1})/\phi(h) = \nu(h)$ 的含义是：$\phi$ 在 $\tau$ 的共轭作用下的「形变」恰为 $\nu$；这等价于说 $\mu\otimes\phi$ 的诱导表示在 $\tau$ 处的交错性恰好复现出 $\sigma\otimes\nu$ 的射影类。$\phi$ 的存在性等价于某个上同调类在 $H^1(\mathbb{Z}_2,\mathbb{Z}_2)$ 中的消没，而 Hasse--Minkowski 定理（对函数域上的二次型给出局部--整体原则）正是用来消去该类的：$\nu$ 作为 $\mathbb{F}_{q^2}(C)/\mathbb{F}_q(C)$ 的二次特征标，在每处局部都能由 $\phi$ 的局部形式实现，整体上才能拼出 $\phi$。Cebotarev 定理在此用于把「$W_F$ 上的恒等式」化归为「Frobenius 共轭类上的恒等式」，因为 $W_F$ 的连续二次特征标由其在 Frobenius 共轭类上的值唯一确定。 来源：原书 p.107 上下文（A）。]

% [背景补充: Hasse--Minkowski 定理（局部--整体原则）在函数域上的形式是：$F=\mathbb{F}_q(C)$ 上两个二次型等价当且仅当它们在所有完备化上等价。原文用它来判定：一个给定的二次特征标 $\nu$ 能否写成 $\phi(\tilde{\tau}h\tilde{\tau}^{-1})/\phi(h)$ 的形式，这可化归为某个二次型（或二元素代数）的不变量在所有局部处相容的整体问题；该原则仅对二次型成立（次数 $\geq 3$ 的形式一般不满足），这正是论证只涉及二次扩张的原因。 译者按（背景补充，无外部依据）。]

因为同一个二次扩张与 $\sigma$ 和 $\sigma'$ 都相伴，我们发现 $\Bun^{\mathcal{O}(D)}_{SL_2}(\mathbb{F}_q)$ 上
% ==== 原文 PDF 第 108 页（印刷页 108）====
关于 $\sigma$ 与 $\sigma'$ 的 Hecke 本征函数空间的维数判据是相同的。因此当 $\#S$ 为奇数时，Hecke 本征函数 $f^D_\sigma$ 与 $f^D_{\sigma'}$ 必须同时消失，其中 $S\subset D$ 是 $D$ 中在 $E$ 内非分裂的点之子集。

现在回忆，对应于 $\Bun^{\mathcal{O}(D)}_{SL_2}$ 上的层 $\mathcal{F}^D_{\sigma,\pm}$ 的函数 $f^D_{\sigma,\pm}$ 由公式 (8.21) 给出：
\begin{equation}\tag{8.24}
f^D_{\sigma,\pm} = \frac{1}{2}\bigl(f^D_\sigma \pm f^D_{\sigma'}\bigr).
\end{equation}
于是我们发现，若 $\#S$ 为奇数，则这两个函数都必须为零。

因此我们得到：当 $\#S$ 为奇数时，$\Bun^{\mathcal{O}(D)}_{SL_2}$ 上的分数 Hecke 本征层 $\mathcal{F}^D_{\sigma,\pm}$ 使得 $\Bun^{\mathcal{O}(D)}_{SL_2}$ 的 $\mathbb{F}_q$-点之集上相应的函数 $f^D_{\sigma,\pm}$ 恒等于 $0$。然而，这并不意味着层本身等于 $0$。这仅仅意味着 $\mathcal{F}^D_{\sigma,\pm}$ 在 $\Bun^{\mathcal{O}(D)}_{SL_2}$ 的 $\mathbb{F}_q$-点处的茎上 Frobenius 的迹等于 $0$。但这并不意味着在 $\mathbb{F}_{q^n}$-点处（$n > 1$）茎上 Frobenius 的迹等于 $0$（若如此则意味着层恒为零，见 [La1]）。事实上，容易看出对一般的 $n$，后者不为零。

在解释这一点之前，我们注意关于 Hecke 本征层与基变换相容性的一个一般事实。对每个 $n > 1$，我们有 $\mathbb{F}_{q^n}$ 上的曲线
\[
C_n = C\times_{\operatorname{Spec}\mathbb{F}_q}\operatorname{Spec}\mathbb{F}_{q^n}.
\]
$C_n$ 上 $G$-丛的模叠 $\Bun_{G,C_n}$ 等价于 $C$ 上 $G$-丛的模叠 $\Bun_G = \Bun_{G,C}$ 的基变换
\[
\Bun_{G,C_n} = \Bun_G\times_{\operatorname{Spec}\mathbb{F}_q}\operatorname{Spec}\mathbb{F}_{q^n}.
\]
设 $\mathcal{E}_n$ 是 $C$ 上 ${}^{L}G$-局部系统 $\mathcal{E}$ 到 $C_n$ 的拉回。几何 Langlands 对应与基变换相容，意思是：以 $\mathcal{E}$ 为本征值的 Hecke 本征层 $\mathcal{F}$ 从 $\Bun_{G,C}$ 到 $\Bun_{G,C_n}$ 的拉回 $\mathcal{F}_n$ 是以 $\mathcal{E}_n$ 为本征值的 Hecke 本征层（见 [La2]）。

现在考虑层 $\mathcal{F}^D_{\sigma,\pm}$ 在 $\Bun^{\mathcal{O}(D)}_{SL_2,C}$ 的 $\mathbb{F}_{q^m}$-点处（$m > 1$）茎上 Frobenius 的迹。作为例子，考察 $\Bun^{\mathcal{O}(D)}_{SL_2,C}$ 的 $\mathbb{F}_{q^2}$-点之集，它等同于 $C_2$ 上 $SL_2$-丛的模叠 $\Bun^{\mathcal{O}(D)}_{SL_2,C_2}$ 的 $\mathbb{F}_{q^2}$-点之集。Hecke 本征层 $\mathcal{F}^D_\sigma$ 到 $\Bun^{\mathcal{O}(D)}_{SL_2,C_2}$ 的拉回是以 $\sigma_2$ 为本征值的 Hecke 本征层 $\mathcal{F}^D_{\sigma_2}$，其中 $\sigma_2$ 是 $\sigma$ 在 $W_{F_2}$ 上的限制，$F_2 = \mathbb{F}_{q^2}(C)$。因此 $\mathcal{F}^D_{\sigma,\pm}$ 的拉回是 $\mathcal{F}^D_{\sigma_2,\pm}$。例如设 $D = [y]$，其中 $y$ 是 $C$ 的 $\mathbb{F}_q$-点，且它在用于定义 $\sigma$ 的 $F$ 的二次扩张 $E$ 中非分裂。则定理 7.1 陈述中出现的集合 $S$ 仅由一个点 $y$ 组成，而按该定理，与 $\mathcal{F}^D_{\sigma,\pm}$ 相联系的 $\Bun^{\mathcal{O}(y)}_{SL_2,C}(\mathbb{F}_q)$ 上的函数 $f^{[y]}_{\sigma,\pm}$ 必须为零。

% [原文疑误: 原书 p.108 此处作 "which is is the same as the set of ..."，重复了 is；译文按单次 is 处理。 来源：原书 p.108 英文原文与译本对照（A）。]

然而，$C_2$ 中对应于 $y$ 的 $\mathbb{F}_{q^2}$-点（我们也把它记为 $y$）在 $F_2$ 的相应二次扩张 $E_2$ 中是分裂的。因此由定理 7.1，$\Bun^{\mathcal{O}(y)}_{SL_2,C_2}(\mathbb{F}_{q^2}) = \Bun^{\mathcal{O}(y)}_{SL_2,C}(\mathbb{F}_{q^2})$ 上的函数 $f^D_{\sigma_2,\pm}$ 不为零。根据
% ==== 原文 PDF 第 109 页（印刷页 109）====
上述与基变换的相容性质，这些函数与 $\Bun^{\mathcal{O}(D)}_{SL_2,C}(\mathbb{F}_{q^2})$ 上对应于层 $\mathcal{F}^{[y]}_{\sigma,\pm}$ 的函数重合。因此层本身不为零！

% [Tutor Note: 易错预警。这一段是全文最容易被误读的推理之一：$f_{\mathcal{F},\mathbb{F}_q}\equiv 0$ 绝不蕴含 $\mathcal{F}=0$。Grothendieck 词典只取了 $\mathbb{F}_q$-点处的迹；要判定层消失，须所有 $n\geq 1$ 的 $\mathbb{F}_{q^n}$-点处迹全为零（[La1] 的判据）。原文的论证是：$\#S$ 为奇数时 $\mathbb{F}_q$-点处对应的函数为零，但基变换到 $C_2$ 后同一闭点变得分裂、判据翻转，$\mathbb{F}_{q^2}$-点处函数非零，故层非零。另需注意 $D=[y]$ 的基变换仍记作 $\mathcal{O}(y)$，但此时的 $y$ 已是 $C_2$ 上的点，故 $\Bun^{\mathcal{O}(y)}_{SL_2,C_2}(\mathbb{F}_{q^2})$ 与 $\Bun^{\mathcal{O}(y)}_{SL_2,C}(\mathbb{F}_{q^2})$ 这两个写法指的是同一个集合。 来源：原书 p.109 上下文（A）。]

这个初等的例子表明：即使与层 $\mathcal{F}^D_{\sigma,\pm}$ 相联系的 $\Bun^{\mathcal{O}(D)}_{SL_2,C}$ 的 $\mathbb{F}_q$-点之集上的函数为零，相应的在 $m > 1$ 的 $\mathbb{F}_{q^m}$-点之集上的函数也不会同时全为零，因此层 $\mathcal{F}^D_{\sigma,\pm}$ 非零。

\subsection{阿贝尔情形}

在 6.6 节我们讨论过分数 Hecke 本征层的其他例子。现在我们在底曲线 $C$ 定义在 $\mathbb{F}_q$ 上的情形重新考察它们。考察 $\Bun_G$ 的 $\mathbb{F}_q$-点之集上相应的函数，并像我们上面在 $G = SL_2$ 的内窥例子中所做的那样（见公式 (8.24)），用通常的 Hecke 本征函数来表达它们，是有启发的。

首先考虑 $G$ 是一维环面的情形。相应的模空间即 Picard 簇 $\mathrm{Pic}$，它分解为连通分支 $\mathrm{Pic}_n$，$n\in \mathbb{Z}$；而对应于 Weil 群 $W_F$ 的一维（$\ell$-进）表示 $\sigma$ 的 Hecke 本征层 $\mathcal{F}_\sigma$ 分解为直和
\begin{equation}\tag{8.25}
\mathcal{F}_\sigma = \bigoplus_{n\in\mathbb{Z}} \mathcal{F}_{\sigma,n},
\end{equation}
其中 $\mathcal{F}_{\sigma,n}$ 支撑在 $\mathrm{Pic}_n$ 上。这是分解 (8.6) 的类比物。设 $f_\sigma$（相应地，$f_{\sigma,n}$）是 $\mathrm{Pic}(\mathbb{F}_q)$（相应地，$\mathrm{Pic}_n(\mathbb{F}_q)$）上对应于 $\mathcal{F}_\sigma$（相应地，$\mathcal{F}_{\sigma,n}$）的函数。则我们有
\[
f_\sigma = \sum_{n\in\mathbb{Z}} f_{\sigma,n}.
\]
这类似于公式 (8.10)。现在我们希望像公式 (8.24) 那样，用（通常的）Hecke 本征函数 $f_{\sigma'}$ 来表达函数 $f_{\sigma,n}$。

这通过一个简单的 Fourier 变换实现。即，对每个非零数 $\gamma\in\mathbb{C}^\times$（下文中我们把 $\overline{\mathbb{Q}}_\ell$ 与 $\mathbb{C}$ 等同），我们通过如下同态的复合来定义 $W_F$ 的一维表示 $\alpha_\gamma$：同态
\begin{equation}\tag{8.26}
\mathrm{res} : W_F \to W_{\mathbb{F}_q} = \mathbb{Z},
\end{equation}
（它由限制到标量 $\mathbb{F}_q\subset F$ 得到）与同态
\[
\mathbb{Z}\to \mathbb{C}^\times, \qquad 1\mapsto \gamma.
\]
现在令 $\sigma_\gamma = \sigma\otimes\alpha_\gamma$ 为 $\sigma$ 经 $\alpha_\gamma$ 的扭转。则我们有如下显然的公式
\begin{equation}\tag{8.27}
f_{\sigma,n} = \frac{1}{2\pi i}\int_{|\gamma|=1} f_{\sigma_\gamma}\,\gamma^{-n-1}\mathrm{d}\gamma,
\end{equation}
它把函数 $f_{\sigma,n}$ 表示为对应于 $\sigma$ 经 $\alpha_\gamma$（$|\gamma|=1$）扭转所得的通常 Hecke 本征函数的积分。

公式 (8.27) 是我们在内窥情形中有过的公式 (8.24) 的类比物，其意义是：两种情形中，满足分数 Hecke 性质的函数（即内窥情形中的 $f_{\sigma,\pm}$ 与阿贝尔情形中的 $f_{\sigma,n}$，$n\in\mathbb{Z}$）都通过通常 Hecke 本征函数的 Fourier 变换来表达。区别在于：第一种
% ==== 原文 PDF 第 110 页（印刷页 110）====
情形中，Fourier 变换是对有限群 $\mathbb{Z}_2$（即内窥同态 $\sigma : W_F \to PGL_2$ 的自同构群）进行的；而第二种情形中，它是对一个连续的自同构群（更确切地说，其紧形式 $U_1$）进行的。这就是为什么 (8.24) 中的有限和会被 (8.27) 中的积分取代。我们将在 9.4 节看到关于更一般的有限自同构群作这类 Fourier 变换的其他例子。

% [Tutor Note: 推导补足。(8.26) 中的 $\mathrm{res}$ 是「次数/赋值」映射：$W_F\to W_{\mathbb{F}_q}=\hat{\mathbb{Z}}$ 的对数版本，即把 $W_F$ 的元素送到其作用的标量域上的次数；在 $\mathbb{Z}$ 上取值的这一支恰好由整体类域论给出，故 $\alpha_\gamma$ 就是「非分歧特征标 $\gamma^{\deg}$」。于是 (8.27) 是对紧群 $U_1 = \{|\gamma| = 1\}$ 的 Fourier 反演公式：把函数按角色 $\gamma\mapsto \gamma^n$ 展开，测度 $\frac{1}{2\pi i}\gamma^{-n-1}\mathrm{d}\gamma$ 正是 $U_1$ 上的标准 Haar 测度（弧长测度）。这与 (8.24) 中对 $\mathbb{Z}_2$ 的离散 Fourier 变换（即 $(f_+\pm f_-)/2$）结构相同，只是群从 $\mathbb{Z}_2$ 换成了 $U_1$，故和换成积分——原文正是借此说明「分数 Hecke 性质 = 对自同构群作 Fourier 变换」这一统一图景。 来源：原书 p.110 上下文（A）。]

以类似的方式，我们可以得到与 6.6 节所讨论的其他类型局部系统相联系的、满足分数 Hecke 性质的函数：对更一般的环面、对自同构群为 ${}^{L}G$ 的中心的局部系统、以及对 Eisenstein 级数。为这些函数以及它们所满足的分数 Hecke 性质找到直接的表示论解释，将是有意义的。

\subsection{Iwahori 情形}

上文我们讨论了模叠 $\Bun^{\mathcal{O}(D)}_{SL_2}$ 上的 Hecke 本征层以及相应的 Hecke 本征函数。然而，在第 3 节中我们关于椭圆曲线所讨论的最详细的 A-膜例子中，我们考虑的是一个略有不同的模空间，它对应分歧 Higgs 丛。在这一情形，相关的模叠是 $\Bun^{\mathcal{O}(D)}_{SL_2,I_p}$，它参数化 $C$ 上行列式为 $\mathcal{O}(D)$ 的秩二向量丛，以及在 $C$ 的一个固定点 $p$ 处的抛物结构（即丛在 $p$ 处纤维中一条直线的选取）。考察 7.3 节所讨论的 L-包的故事在这一情形如何展开，是有启发的。

设 $C$ 再次定义在 $\mathbb{F}_q$ 上。则 $\Bun^{\mathcal{O}(D)}_{SL_2,I_p}$ 的 $\mathbb{F}_q$-点之集同构于双商
\begin{equation}\tag{8.28}
SL_2(F)\backslash SL_2(\mathbb{A}_F)\Big/\left(\prod_{x\neq p} K_x\times I_p\right).
\end{equation}
这里 $I_p = K'_p\cap K''_p$ 是 $SL_2(F_p)$ 的 Iwahori 子群，且对 $x\in D$ 有 $K_x = K''_x$，否则 $K_x = K'_x$。设 $p$ 是 $C$ 的非分裂点，这是相对于与某个非分歧同态 $\sigma : W_F \to O_2$ 相伴的非分歧覆盖 $C'\to C$ 而言的。则对应于 $p$ 与按上述方式构造的同态 $\sigma : W_F \to PGL_2$ 的局部 L-包由两个不可约表示 $\pi'_p$ 与 $\pi''_p$ 组成，但现在 $(\pi'_p)^{I_p}$ 与 $(\pi''_p)^{I_p}$ 都是一维的。

固定局部因子 $\pi_x$（$x\neq p$）。则我们有两个互不同构的 $SL_2(\mathbb{A}_F)$ 的不可约表示
\[
\bigotimes_{x\neq p}\pi_x\otimes \pi'_p \qquad\text{and}\qquad \bigotimes_{x\neq p}\pi_x\otimes \pi''_p.
\]
按定理 7.1，其中只有一个是自守的；也就是说，它可以实现为 $SL_2(F)\backslash SL_2(\mathbb{A}_F)$ 上适当函数空间的一个组分。然而，它们关于子群
\[
\prod_{x\neq p} K_x\times I_p
\]
的不变空间都是一维的。因此无论其中哪一个在自守意义下出现，我们都会在双商
% ==== 原文 PDF 第 111 页（印刷页 111）====
(8.28) 上得到一维的 Hecke 本征函数空间。于是，与正则 Hecke 本征层相联系的 $\Bun^{\mathcal{O}(D)}_{SL_2,I_p}$ 的 $\mathbb{F}_q$-点之集上的函数将不为零。如上所述，我们随后得到与分数 Hecke 本征层相联系的函数 $f^D_\pm$ 在这一情形也不为零。这构成了双商 (7.12) 与 (8.28) 之间的一个重要差别。

% [Tutor Note: 物理图像。Iwahori 情形之所以消除 L-包的消失现象，关键在于 $I_p = K'_p\cap K''_p$ 比 $K'_p$、$K''_p$ 都小，故 $I_p$-不变空间比 $K'_p$-不变空间、$K''_p$-不变空间都大：局部 L-包的两个成员在 $I_p$-不变空间中都有非零向量，且在新情形二者都是一维的。于是「哪一个成员自守」不再改变不变空间的维数，非 Iwahori 的双商 (7.12) 中那种「某个成员的不变空间为零、函数被迫消失」的现象也就不再出现。物理上，$p$ 处取抛物结构对应 Higgs 丛在该点分歧（在原书 §3 的椭圆曲线例子中即驯顺分歧点），它改变了 Hitchin 纤维在该处的局部模型，从而使分数本征膜不再被迫成对消失。 来源：原书 p.111 上下文（A）。]
